\section{Unified Finite Converse}\label{sec:converse}
%==============================================================================

\leanmetapending{\LHrng{OBS}{1}{2}}
\begin{theorem}[Pair-injectivity characterization]\label{thm:pair-injective}
Let $X$ range over a finite latent alphabet, let $Y$ be the deterministic observation transcript, and let $T$ be a finite auxiliary tag. Zero-error recovery from $(Y,T)$ is possible if and only if the pair map $x \mapsto (Y(x),T(x))$ is injective on the ambiguity class under consideration.
\end{theorem}

\begin{proof}
Let $\phi(x):=(Y(x),T(x))$.

\emph{Necessity.} If $\phi$ is not injective, then there exist distinct latent states $x_1\neq x_2$ with $\phi(x_1)=\phi(x_2)$. Any deterministic decoder $D$ therefore receives the same input on both states and must satisfy $D(\phi(x_1))=D(\phi(x_2))$. Since $x_1\neq x_2$, the common output cannot equal both latent states, so $D$ errs on at least one of them.

\emph{Sufficiency.} If $\phi$ is injective, define a decoder on the image of $\phi$ by $D(y,t)=\phi^{-1}(y,t)$, with arbitrary extension off the image. Injectivity makes $\phi^{-1}$ single-valued on the image, and for every latent state $x$ we then have $D(Y(x),T(x))=x$. Thus zero-error recovery is possible.
\end{proof}

\leanmetapending{\LH{OBS5}}
\begin{theorem}[Confusability formulation]\label{thm:confusability-converse}
Fix a deterministic observation transcript and an $L$-bit side-information tag. If $K$ latent states induce the same observation transcript, then zero-error decoding on that ambiguity class requires $K$ distinct tag outcomes. Equivalently, a $K$-way ambiguity class forms a confusability clique whose size cannot exceed the available tag alphabet.
\end{theorem}

\begin{proof}
Apply Theorem~\ref{thm:pair-injective} on an ambiguity class on which the observation is constant. Then injectivity of $(Y,T)$ reduces to injectivity of the tag coordinate alone on that class. Since an $L$-bit tag provides at most $2^L$ outcomes, the clique size cannot exceed $2^L$.
\end{proof}

\leanmetapending{\LH{CIA3}}
\begin{corollary}[Counting converse]\label{thm:fano-converse}
Theorem~\ref{thm:confusability-converse} specialized to a single $K$-way ambiguity class gives $K \le 2^L$, i.e., at least $\log_2 K$ bits of side information.
\end{corollary}

\medskip
\noindent\textbf{Proposition III.4 (Equivalent formulations).} The pair-injectivity obstruction of Theorem~\ref{thm:pair-injective} admits the following equivalent formulations on a finite source observed through a deterministic transcript with a finite auxiliary tag:
\begin{itemize}
\tightlist
\item \emph{Global budget:} $K \le O\,T$ for latent size $K$, observation alphabet size $O$, and tag alphabet size $T$.\leanmeta{\LH{OBS4}}
\item \emph{Fiber injectivity:} $|\mathcal F_y| \le |\mathcal T|$ for each observation fiber $\mathcal F_y=\{x : Y(x)=y\}$.\leanmeta{\LH{OBS3}}
\item \emph{Conditional entropy:} $H(X \mid Y) = \log_2 K$ on a $K$-way class with constant observation, and $H(T \mid Y) \ge \log_2 K$.\leanmeta{\LH{PRB47}, \LH{PRB55}}
\item \emph{Decoder output:} $H(\hat X) \le H(Y,T) \le H(X)$, with nonnegative gap to alphabet ceilings.\leanmeta{\LHrng{PRB}{73}{74}, \LHrng{PRB}{82}{83}, \LH{PRB90}, \LHrng{PRB}{93}{95}}
\item \emph{Mass-weighted clique:} $\sum_{x\in S} p(x)\log_2 \frac{1}{p(x)} \le h_2(p_S) + p_S \log_2 |\mathcal T|$ for a success-set clique $S$.\leanmeta{\LH{PRB55}}
\item \emph{Data processing:} $H(\kappa(Y,T)) \le H(Y,T)$ for any deterministic coarsening $\kappa$.\leanmeta{\LH{PRB68}, \LH{PRB81}, \LH{PRB83}}
\end{itemize}

\emph{Proof sketch.} Each formulation follows from Theorem~\ref{thm:pair-injective} by restricting the injectivity condition to the relevant domain (fiber, clique, full source) and applying standard entropy identities. The confusability, counting, conditional-entropy, decoder-output, and finite-gap statements are therefore different normal forms of the same obstruction: a surviving $K$-way ambiguity class requires a budget of at least $\log_2 K$ bits to be resolved exactly.\leanmeta{\LH{OBS1}, \LH{OBS5}, \LH{CIA3}}

\subsection{Finite-error extension}\label{sec:finite-error-extension}

The main body of the paper studies zero error. The same deterministic finite model also supports a budgeted finite-error extension, which should be read as a deterministic finite analogue of the classical Fano line of argument rather than as a separate asymptotic coding theorem~\cite{fano1961transmission,witsenhausen1975conditional,cover2006elements,kelly2011improved}.

\leanmetapending{\LH{PRB85}, \LH{PRB87}}
\begin{proposition}[Budgeted finite-error converse]\label{thm:finite-error-budgeted}
Let $P_e$ denote the decoder error probability on a finite source of size $K$, observed through a deterministic transcript alphabet of size $O$ together with a tag alphabet of size $T$. Then the decoded-output entropy obeys
\[
H(\hat X)
\;\le\;
h_2(P_e) + (1-P_e)\log_2 (OT) + P_e \log_2 (K-1).
\]
\end{proposition}

\begin{proof}
Partition the source into success and failure events. On the success branch, the decoded output is constrained by the observation-tag budget and contributes at most $\log_2(OT)$. On the failure branch, the decoder can still output at most one of the remaining $K-1$ alternatives. The Bernoulli split between the two branches contributes the binary-entropy term. Equivalently,
\[
H(\hat X)=H(\hat X \mid \text{success/failure}) + H(\text{success/failure})
\]
is bounded by combining the support bound on each branch with the binary entropy of the branch variable itself.
\end{proof}

The zero-error theorems are the $P_e=0$ boundary of this inequality. In that regime the success branch occupies all probability mass, the Bernoulli term vanishes, the failure contribution disappears, and the budget collapses back to the unified finite converse above.
